\section{Các nghiên cứu liên quan}\label{sec:related}

\subsection{Mô hình người học và truy vết tri thức}

Truy vết tri thức ước lượng tri thức đang thay đổi của người học từ chuỗi câu trả lời. Truy vết tri thức Bayes mô hình hóa mỗi thành phần tri thức như một trạng thái nhị phân ẩn \citep{corbett1994bkt}; các mô hình logistic như mô hình nhân tố cộng (additive factor model, AFM) và phân tích nhân tố thành tích (performance factors analysis, PFA) kết hợp năng lực người học, độ khó của bài và số lần luyện tập theo từng thành phần tri thức \citep{cen2006lfa,pavlik2009pfa,pelanek2017learner}, theo nghĩa thành phần tri thức mà \citet{koedinger2012kli} mô tả. Các mô hình hồi quy và mô hình dựa trên cơ chế chú ý, khởi đầu với DKT \citep{piech2015dkt} rồi SAKT và AKT \citep{pandey2019sakt,ghosh2020akt}, học động học trực tiếp từ dữ liệu \citep[xem tổng quan của][]{abdelrahman2023survey}. Nhiều so sánh cho thấy các mô hình logistic được đặc tả tốt đạt hoặc tiệm cận độ chính xác dự báo của mô hình sâu \citep{khajah2016deep,wilson2016basics,gervet2020deep}. Dự báo từng bước của DKT cũng có thể dịch chuyển theo hướng phi lý, chẳng hạn hạ mức thành thạo dự báo của một kỹ năng sau một câu trả lời đúng \citep{yeung2018dktplus}, dẫn tới các ràng buộc nhất quán và đơn điệu \citep{zhang2023cmkt}. Với lập trình, KT đã được mở rộng bằng đặc trưng mã nguồn \citep{shi2022codedkt}, và hỗ trợ điều hướng thích ứng từ lâu đã dẫn người học tới các bài lập trình có độ khó phù hợp \citep{hsiao2010guiding,brusilovsky2007user}. Dù vậy, việc đánh giá mô hình người học vẫn bị chi phối bởi độ chính xác dự báo trên dữ liệu giữ lại, mặc dù đã có những khung đánh giá mô hình qua hệ quả của nó đối với các quyết định dạy học \citep{gonzalezbrenes2015leopard,huang2015multifaceted}.

\subsection{Học tăng cường cho sắp xếp trình tự dạy học và gợi ý bài tập}

Sắp xếp trình tự dạy học bằng RL có lịch sử lâu dài. Các chính sách rút ra từ dữ liệu đã được thử với học sinh trong hệ thống dạy kèm thông minh \citep{chi2011umuai,shen2018jedm,zhou2022hrl}, trong khóa học trực tuyến \citep{bassen2020chi} và trong một hệ dạy toán dạng câu chuyện \citep{ruan2024rltutor}; các hướng liên quan dùng lập kế hoạch POMDP \citep{rafferty2016pomdp}, bandit nhiều cần gạt \citep{clement2015bandits} và mô hình dạy học xấp xỉ tối ưu \citep{whitehill2018teaching}. Các tổng quan kết luận rằng thành công là có thật nhưng không đồng đều: \citet{doroudi2019reward} nhận thấy chính sách rút ra từ mô hình người học thuần dữ liệu thường không vượt các đường cơ sở hợp lý khi thử với học sinh, và thành công dễ đạt hơn khi chính sách dựa trên mô hình có cơ sở lý thuyết; các tổng quan gần đây hơn chỉ ra các vấn đề phương pháp luận và nhu cầu đánh giá quy mô lớn hơn với người học thật \citep{singla2021rleducation,riedmann2025rlreview}. Song song, cộng đồng hệ gợi ý đã xây dựng các bộ gợi ý bài tập và lộ trình học bằng RL sâu, dựa vào một mô hình người học đã fit — thường là mô hình KT — để cung cấp trạng thái, phần thưởng hoặc thước đo việc học: gợi ý bài tập đa mục tiêu \citep{huang2019dre}, lộ trình học có nhận biết cấu trúc nhận thức \citep{liu2019cseal}, gợi ý dựa trên DKT có nhận biết khái niệm \citep{ai2019cdkt}, các bộ lập lộ trình dạng tập-thành-chuỗi và phân cấp \citep{chen2023src,li2023gehrl}, chưng cất trạng thái tri thức \citep{li2024pksd}, và ExRec, định nghĩa phần thưởng và điểm nhiệm vụ qua các trạng thái KT đã hiệu chỉnh \citep{ozyurt2025exrec}. Benchmark UniER hợp nhất các bộ gợi ý mức bài và mức lộ trình dưới một chỉ số tiến bộ nhận thức có trọng số rút từ KT \citep{cheng2026unier}, và các hệ lộ trình học bằng RL sâu cũng đã xuất hiện trên tạp chí này \citep{ruan2025adaptive}. Các công trình đó cải tiến biểu diễn và tối ưu hóa; câu hỏi của chúng tôi là bao nhiêu phần của một tiến bộ đo bằng KT còn đứng vững khi thế giới sinh ra kết quả không phải là mô hình đo kết quả.

\subsection{Từ độ chính xác dự báo đến quyết định dạy học}

Việc độ chính xác dự báo không quyết định được các lựa chọn dạy học là một phát hiện lặp lại. Với một tỉ lệ đáng kể học sinh, tham số truy vết tri thức chung của quần thể hàm ý lượng luyện tập khoảng gấp đôi so với tham số fit riêng cho từng người \citep{lee2012individualizing}; các mô hình người học có độ chính xác dự báo tương tự có thể hàm ý lượng luyện tập khác nhau đáng kể \citep{rollinson2015predictive}; tham số truy vết tri thức có thể suy biến về mặt ngữ nghĩa dù mô hình fit tốt \citep{doroudi2017identifiability}; và mô hình đặc tả sai có thể khiến quyết định về mức thành thạo thiếu công bằng giữa người học \citep{doroudi2019fairer}. \citet{gonzalezbrenes2015leopard} đề xuất đánh giá mô hình người học qua các quyết định dạy học mà nó tạo ra. Ma trận đánh giá bền vững \citep{doroudi2017rem} đánh giá mỗi chính sách ứng viên dưới nhiều mô hình người học và cảnh báo việc tin vào một chính sách chỉ tốt dưới chính mô hình đã sinh ra nó. Các bộ ước lượng off-policy dựa trên lấy mẫu theo trọng số quan trọng (importance sampling) đánh giá chính sách từ log mà không cần bộ mô phỏng \citep{mandel2014offline,thomas2016dataefficient}, nhưng chúng đòi hỏi xác suất hành động được ghi lại, điều mà việc tự luyện tập trên hệ thống chấm bài không cung cấp, nên mô phỏng vẫn là cách mặc định để đánh giá bộ gợi ý bài tập. Trong học máy nói chung, các bộ dự báo có độ chính xác tương đương trên dữ liệu giữ lại có thể hành xử rất khác nhau khi dùng cho tác vụ hạ nguồn \citep{breiman2001two,damour2022underspecification}.

\subsection{Tối ưu trên mô hình học từ dữ liệu và phần thưởng đại diện}

Trong RL dựa trên mô hình, một chính sách tối ưu trong mô hình động học học từ dữ liệu sẽ khai thác sai số của mô hình, nhất là ở vùng thiếu dữ liệu; các kỹ thuật tổ hợp mô hình (ensemble), mô phỏng ngắn và phạt độ bất định được đưa ra để hạn chế điều này \citep{kurutach2018metrpo,janner2019trust,yu2020mopo,moerland2023mbrl,levine2020offline}. Tương tự, tối ưu một mục tiêu đại diện có thể làm tăng mục tiêu đại diện trong khi mục tiêu thật đứng yên — hiện tượng lách hay chơi gian phần thưởng (reward hacking/gaming) \citep{amodei2016concrete,skalse2022reward} — và việc tối ưu quá mức một mô hình phần thưởng học từ dữ liệu tuân theo các quy luật co giãn dự đoán được \citep{gao2023scaling}. Bộ mô phỏng của hệ gợi ý thừa hưởng thiên lệch của log dùng để fit nó \citep{huang2020keeping}, và đánh giá ngoại tuyến là một trở ngại đã được thừa nhận khi triển khai RL trong thế giới thực \citep{dulacarnold2021challenges}. Bộ gợi ý KT--RL nằm ở giao điểm của các vấn đề này vì một mô hình học từ dữ liệu đóng cùng lúc vai trò động học, phần thưởng và thước đo đánh giá. Từ dịch tễ học, chúng tôi mượn khái niệm đối chứng âm: một phân tích mà hiệu ứng thật được biết chắc là không tồn tại, nên mọi hiệu ứng ước lượng được đều bộc lộ sai lệch \citep{lipsitch2010negative,arnold2016negative}.

Công trình này nối các mạch nghiên cứu trên cho bài toán gợi ý bài tập lập trình. Chúng tôi fit sáu mô hình người học có chất lượng dự báo ngang nhau trên cùng dữ liệu thật và cho thấy sự bất đồng giữa chúng lan truyền sang thứ hạng chính sách như thế nào, đồng thời bổ sung hai công cụ mà các nghiên cứu về đánh giá bền vững chưa dùng: một thế giới đối chứng âm không có học, trong đó tiến bộ thật của mọi chính sách đúng bằng 0, và một phép phân rã chính xác tiến bộ mà tracker báo cáo thành độ trôi và phần bất ngờ, kèm các kết quả hình thức đồng nhất tiến bộ trong mô phỏng với độ trôi tích lũy.
